% TeX-ready mathematical fragment. Requires amsmath, amssymb, and booktabs.
\section{A short local majorant for nine positive variables}

A separate reduction gives two explicit nine-positive-variable equations:
\[
 (442,\,1\,520\,186),\qquad (456,\,1\,376\,378),
\]
where the entries are live binary arithmetic operations and exact total degree.
The second improves both measures over the reviewed
$475$-operation, degree-$1\,379\,730$ equation with the same original31 tuple.
The first is a tradeoff against the $433$-operation,
degree-$1\,628\,768$ compact54-sharp equation.  Neither changes the
nine-variable bound.  The new construction is a short local Laurent
majorant combined with a permissible successor coordinate.  No global
minimality or historical priority is asserted.

\subsection{The local lemma}

Let $A\ge2$, $C\ge2$, $i\ge1$ be integers and put
\begin{align*}
 J&=4A-5,& \delta&=A^2-1,&D&=1+\delta C^2,\\
 e&=iJC^2,&F&=1+\delta De^2,&
 T&=iC^2\bigl[AJ(2D-C^2)+3C^2-4\bigr].
\end{align*}
Then
\begin{equation}\label{eq:short-local-majorant}
                  T>2\sqrt{DF}>0.
\end{equation}
To prove this, set $t=C^2\ge4$. Exact expansion gives
\[
 T^2-4DF=i^2t^2Q(A,t)-4D,
\]
where
\begin{align*}
 Q(A,t)={}&20t(3t-4)A^3-t(83t-116)A^2\\
          &-10(7t^2-14t+8)A+109t^2-224t+116.
\end{align*}
Writing $A=x+2$, $t=y+4$, the coefficients of
$Q(x+2,y+4)-(x+2)^2$ are
\[
\begin{array}{c|rrr}
 &y^0&y^1&y^2\\ \hline
 x^0&1344&816&117\\
 x^1&3580&2188&318\\
 x^2&2975&1852&277\\
 x^3&640&400&60
\end{array}
\]
and are all positive. Consequently $Q(A,t)>A^2$ and
\[
 T^2-4DF>t^2A^2-4\bigl[(A^2-1)t+1\bigr]
       =A^2t(t-4)+4t-4\ge12.
\]
The defining expression gives $T>0$, since
$2D-t=2+(2A^2-3)t>0$, $J>0$, and $3t-4>0$.
This proves~\eqref{eq:short-local-majorant}.

The boundary condition is essential: at $(A,C,i)=(3,1,1)$ one obtains
$T^2-4DF=-308$. It is therefore supplied before the final equation is
assumed, by replacing its carrier input with $C=c+1$ for a positive
unknown $c$. Canonical Pell witnesses already have $C>1$, so this
restriction preserves existential completeness.

\subsection{The complete polynomial}

The nine positive unknowns are $b,g,h,s,w,c,i,j,n$. The two inherited
source prefixes define $R,N,\operatorname{cap}$ from their complete
literal coefficient tables. The delivered package includes both tables,
all four-port numerals, all fixed circuit leaves, and the exact prefix
maps. The whole compact54-sharp and original31 tuples are kept separate.
Both source prefixes prove $R,N\ge8$ and $b\le R,N$ on the ambient
positive domain.

Set $C=c+1$ and retain
\begin{align*}
 U&=N^2w,&Y&=N^2s,&M&=RY,&P&=2M^2U,\\
 K&=R+1+h(P-1),&A&=M(U+1),&B_1&=2R+1,&W_1&=bw,\\
 J&=4A-5,&Q&=3W_1C-2(W_1^2-1),\\
 a_0&=1+(P^2-1)K^2,&V_0&=PK,\\
 m&=K^2(4\operatorname{cap}-g)
       -16\operatorname{cap}(C-KY)^2,\\
 \Omega&=Jm,&E_0&=J(n+W_1^2)+Q.
\end{align*}
For $D,F,e$ as above define
\begin{align*}
 \Lambda&=CD+1-A,&G&=A+F\Lambda,\\
 Z&=B_1+(j-1)C,&H&=C+FZ,&I&=1+(G^2-1)H^2,\\
 V_1&=TGH,&a_1&=4DFI,&\Theta&=E_0-\Omega(V_0+V_1).
\end{align*}
The equation is the ordinary integer polynomial
\begin{equation}\label{eq:local-majorant-nine}
 \bigl[\Theta^2+\Omega^2(a_0-a_1)\bigr]^2
       -4\Omega^2a_0\Theta^2=0.
\end{equation}
All named quantities are polynomial abbreviations, and add no unknowns.
The circuit counts charge the successor $C=c+1$, scalar products, and
all squares. There are no divisions or variable exponents.

\subsection{The unbounded equivalence proof}

The source bounds give $A\ge33280$. Also $D,F,G,H>0$, $G>A$, and
$H>C\ge2$. Therefore $I<(GH)^2$, and the new lemma gives
$V_1>\sqrt{a_1}$. Independently,
$V_0^2-a_0=K^2-1>0$, so $V_0>\sqrt{a_0}$.
These inequalities hold before a zero is assumed.

At a zero, $E_0=Jn+(J-2)W_1^2+3W_1C+2>0$ excludes $\Omega=0$.
Factoring the real two-root norm yields signed roots
$r_0^2=a_0$, $r_1^2=a_1$ with
\[
 E_0=\Omega(V_0+V_1+r_0+r_1).
\]
Both strict majorants make the bracket positive for all sign choices.
Hence $m>0$. Only now do the inherited margin estimates give
$|C-KY|<K/2$, $C>K$, and $C>B_1+W_1$.
Moreover $A^2>P$, $D>C^2$, $F\ge A^2$, and $I>A^2$, whence
\[
 DFI>C^2A^4>K^2P^2>a_0.
\]
Thus $x=r_0+r_1=\Theta/\Omega$ is nonzero and rational. The isolation
identity $2xr_0=x^2+a_0-a_1$ makes both roots rational, hence integral.
The rational number $r_1/2$ has integer square $DFI$, so it too is
integral. The norm relation now gives $J\mid Q$.

The retained saturated Pell theorem recovers $C=\psi_A(B_1)$ from the
square $DFI$; all its range hypotheses have now been established.
Together with $a_0$ square, $m>0$, and $J\mid Q$, these are the exact
inherited ratio conditions, which recover the encoded source zero.

Conversely, a source zero supplies canonical ratio witnesses and compatible
saturated-Pell auxiliaries with $W_1^2+Q/J<2D$. Its carrier satisfies
$C=\psi_A(B_1)>1$, so $c=C-1$ is positive. Choose positive roots
$u_0=\sqrt{a_0}$ and $u_1=2\sqrt{DFI}$, and define
\[
 n=m(V_0+V_1+u_0+u_1)-W_1^2-Q/J.
\]
This is an integer. Since $F>D$, the lemma gives $V_1>2D$, so $n>0$.
Then $\Theta=\Omega(u_0+u_1)$ and
\eqref{eq:local-majorant-nine} holds. This proves both existential
directions. Negation and the parents' retained source/refutation
calibration give the Con(ZFC) biconditional, with neither consistency nor
inconsistency assumed. No new Lean, Rocq, or object-ZFC proof term is claimed.

\subsection{Cancellation-resolved exact degree}

Let $r=\deg R$, $\nu=\deg N$, $a=r+4\nu+2$, and $W=116$.
The unchanged source has an isolated positive leading kernel coefficient.
The shared degrees are
\begin{align*}
 \deg\Omega&=5a+2W+3,&\deg V_0&=3a+r+1,\\
 \deg E_0&=a+4,&\deg\delta_0&=3a+r+2,
\end{align*}
where $\delta_0=V_0^2-a_0$.
In weights $(a,1,1,r)$ on $(A,C,i,Z)$, the unique leading monomials are
\begin{align*}
 [T]_{\rm top}&=8A^4C^4i,\\
 [T^2-4DF]_{\rm top}&=60A^3C^8i^2,\\
 [GH]_{\rm top}&=256A^{14}C^{15}i^4Z,\\
 [V_1]_{\rm top}&=2048A^{18}C^{19}i^5Z,\\
 [\delta_1]_{\rm top}&=3932160A^{31}C^{38}i^{10}Z^2,
\end{align*}
where $\delta_1=V_1^2-a_1$. The last line follows from
\[
 \delta_1=(T^2-4DF)(GH)^2+4DF(H^2-1).
\]
Thus $\deg V_1=18a+r+24$ and $\deg\delta_1=31a+2r+48$.
Set
\[
 B_* = E_0^2-2E_0\Omega(V_0+V_1)+\Omega^2(\delta_0+\delta_1).
\]
The exact outer identity is
\[
 F_9=B_*^2+4\Omega^2V_0V_1B_*
       +4\Omega^4(\delta_0V_1^2+a_0\delta_1).
\]
Its four summand degrees are
\[
 82a+4r+8W+108,\quad72a+4r+8W+85,\quad
 59a+3r+8W+62,\quad57a+4r+8W+62.
\]
The first is strictly greatest and has positive nonzero leading coefficient.
The exact degree is therefore
\[
                        86r+328\nu+8W+272.
\]
For $(r,\nu)=(8799,2324)$ this is $1\,520\,186$;
for $(r,\nu)=(7943,2110)$ it is $1\,376\,378$.
The independent checker also computes nonzero specialized leading
coefficients modulo three primes using the actual fixed numerals.
Those coefficient witnesses supplement the universal identity and
unique-leader argument.

\subsection{Validation and comparison}

The independent standard-library verifier imports no builder. It verifies
all fixed numerals and their recipes, complete prefix maps, every new
register, local positivity and norm identities, topology, liveness, exact
operation counts, degree attainment, and full deterministic regeneration.
Normal and optimized executions agree. It rejects the invalid $C=1$
boundary and an altered majorant gate.

The operation breakdown is
\[
\begin{array}{c|rrrr}
 &+&-&\times&\text{total}\\ \hline
 \text{compact2082}&159&22&261&442\\
 \text{H1868}&155&22&279&456
\end{array}
\]
with exactly nine positive variables in each case. The whole coefficient
bindings and inherited ordinary proofs remain explicit. Finite circuit
checks are not substituted for the unbounded native/compiler/refutation
calibration, and no global simplicity optimum is claimed.
